%% ===========================================================================
%%  1. Introduction
%%  Quantifying with Confidence: Conformal Prediction for Reliable Musical
%%  Style Similarity Ranking
%% ===========================================================================

\section{Introduction}

Musical style similarity assessment is a fundamental problem in music
information retrieval (MIR) and a critical component in the evaluation of
AI-generated music. Whether curating playlists, recommending tracks, or
filtering outputs from generative models such as MusicGen~\cite{copet2023musicgen}
and AudioLDM~2~\cite{liu2023audioldm2}, practitioners need to determine how closely
a piece of music aligns with a target stylistic reference. The reliability of
these similarity judgments directly impacts the quality of downstream
applications, yet the methods used to produce them have largely remained
point-estimate heuristics with no principled uncertainty quantification.

Recent advances in self-supervised learning (SSL) have yielded powerful audio
representation models capable of capturing high-level musical attributes
including timbre, harmony, rhythm, and genre. Among these, the Music
undERstanding model with large-scale self-supervised Training
(MERT)~\cite{li2023mert}---a 330M-parameter Transformer pre-trained on 160,000
hours of music data---has demonstrated state-of-the-art performance across a
wide range of MIR tasks. By extracting embeddings from MERT's deep hidden
representations, one can compute pairwise cosine similarity between tracks
and obtain a single scalar score reflecting their stylistic proximity.
However, this similarity score is a point estimate: it conveys no information
about its own reliability, does not account for the inherent ambiguity of
musical style boundaries, and offers no statistical guarantees. In
safety-critical or curation-sensitive settings---such as filtering
AI-generated music for stylistic compliance or constructing evaluation
benchmarks for generative models---the absence of uncertainty quantification
constitutes a significant limitation.

Conformal Prediction (CP)~\cite{vovk2005algorithmic,shafer2008tutorial}
provides a distribution-free, finite-sample framework for constructing
prediction sets with rigorous coverage guarantees. Given a user-specified
significance level $\alpha$, CP produces a prediction set $\Gamma(x)$ that
is guaranteed to contain the correct label with probability at least
$1 - \alpha$, under only the assumption that calibration and test data are
exchangeable. Unlike Bayesian methods, CP makes no parametric assumptions
about the underlying distribution; unlike bootstrap approaches, it offers
provable finite-sample validity rather than asymptotic approximations.
These properties make CP particularly well-suited to subjective perceptual
tasks such as musical style assessment, where the score distribution is
inherently heavy-tailed and resists simple parametric modeling.

In this paper, we propose the first integration of conformal prediction with
self-supervised audio embeddings for reliable musical style similarity
ranking. Our framework extracts 1024-dimensional embeddings from
MERT-v1-330M for each audio track, constructs per-genre prototype centroids
from a reference set, and defines a nonconformity score based on the cosine
distance between a query embedding and each genre centroid. Using a held-out
calibration set, we compute empirical $p$-values and construct prediction
sets with distribution-free coverage guarantees. We evaluate our method on
the GTZAN dataset~\cite{tzanetakis2002gtzan} (999 tracks across 10 genres
after excluding one corrupted file), comparing against bootstrap resampling
and Gaussian approximation baselines under a stratified 5-fold cross-validation
protocol with a 60/20/20 (reference/calibration/test) split.

We make the following contributions:

\begin{enumerate}
    \item We propose the first conformal prediction framework for musical style
    similarity assessment, bridging the gap between self-supervised audio
    representation learning and statistically rigorous uncertainty
    quantification (Section~\ref{sec:method}).

    \item We demonstrate empirically that our CP approach achieves empirical
    coverage tightly matching nominal levels: $0.993 \pm 0.010$ at
    $\alpha = 0.01$ (target $0.99$), $0.955 \pm 0.021$ at $\alpha = 0.05$
    (target $0.95$), and $0.910 \pm 0.028$ at $\alpha = 0.10$
    (target $0.90$). In contrast, the Gaussian approximation baseline
    yields only $0.965$ coverage at $\alpha = 0.01$, confirming that the
    normality assumption fails in the heavy-tailed regime characteristic
    of perceptual similarity tasks (Section~\ref{sec:experiments}).

    \item We provide a systematic comparison of CP against bootstrap
    resampling and Gaussian approximation baselines, highlighting that
    only CP offers a distribution-free, finite-sample coverage guarantee---a
    critical property when the calibration set is small or the underlying
    distribution is non-Gaussian (Section~\ref{sec:experiments}).
\end{enumerate}

The remainder of this paper is organized as follows. Section~\ref{sec:related}
reviews related work on self-supervised music representation, conformal
prediction, and uncertainty quantification in MIR. Section~\ref{sec:method}
presents our CP framework, detailing the MERT embedding pipeline, genre
centroid construction, nonconformity score design, and prediction set
formation. Section~\ref{sec:experiments} describes our experimental setup,
compares CP against baselines, and analyzes coverage and set size results.
Section~\ref{sec:discussion} discusses the implications of our findings,
including the theoretical advantages of CP in perceptual tasks and
limitations of the current approach. Section~\ref{sec:conclusion} concludes
with a summary of contributions and directions for future work.
